\subsection{Pre-Effect Memory Evaluation}
For a candidate memory operation, the processor validates the instruction and operands, resolves its condition or
predicate, computes the effective address, and completes segmentation, translation, permission, and alignment
validation before evaluating triggers. A winning trigger is reported before the memory or device access and before
any architectural result commits. Consequently, earlier synchronous faults take priority; a matching load performs
no read, a matching store is not visible, a matching atomic performs no read-modify-write action, and a matching MMIO
operation does not access the device. A bus or data fault can still occur after event return when the operation is
retried.

Explicit and implicit instruction data accesses participate. Instruction fetch, page-table walks, event-entry frame
accesses, prefetch hints, and cache-maintenance accesses do not. False-conditioned operations and inactive vector
lanes do not participate. REP and REPcc evaluate before each iteration's memory effect and preserve completed prior
iterations. Gather and scatter evaluate before the next active lane access and preserve completed lane and predicate
progress.
